% Options for packages loaded elsewhere
\PassOptionsToPackage{unicode}{hyperref}
\PassOptionsToPackage{hyphens}{url}
\documentclass[
]{article}
\usepackage{xcolor}
\usepackage{amsmath,amssymb}
\setcounter{secnumdepth}{-\maxdimen} % remove section numbering
\usepackage{iftex}
\ifPDFTeX
  \usepackage[T1]{fontenc}
  \usepackage[utf8]{inputenc}
  \usepackage{textcomp} % provide euro and other symbols
\else % if luatex or xetex
  \usepackage{unicode-math} % this also loads fontspec
  \defaultfontfeatures{Scale=MatchLowercase}
  \defaultfontfeatures[\rmfamily]{Ligatures=TeX,Scale=1}
\fi
\usepackage{lmodern}
\ifPDFTeX\else
  % xetex/luatex font selection
\fi
% Use upquote if available, for straight quotes in verbatim environments
\IfFileExists{upquote.sty}{\usepackage{upquote}}{}
\IfFileExists{microtype.sty}{% use microtype if available
  \usepackage[]{microtype}
  \UseMicrotypeSet[protrusion]{basicmath} % disable protrusion for tt fonts
}{}
\makeatletter
\@ifundefined{KOMAClassName}{% if non-KOMA class
  \IfFileExists{parskip.sty}{%
    \usepackage{parskip}
  }{% else
    \setlength{\parindent}{0pt}
    \setlength{\parskip}{6pt plus 2pt minus 1pt}}
}{% if KOMA class
  \KOMAoptions{parskip=half}}
\makeatother
\usepackage{longtable,booktabs,array}
\usepackage{caption}
\captionsetup[table]{skip=6pt}
\newcounter{none} % for unnumbered tables
\usepackage{calc} % for calculating minipage widths
% Correct order of tables after \paragraph or \subparagraph
\usepackage{etoolbox}
\makeatletter
\patchcmd\longtable{\par}{\if@noskipsec\mbox{}\fi\par}{}{}
\makeatother
% Allow footnotes in longtable head/foot
\IfFileExists{footnotehyper.sty}{\usepackage{footnotehyper}}{\usepackage{footnote}}
\makesavenoteenv{longtable}
\usepackage{graphicx}
\makeatletter
\newsavebox\pandoc@box
\newcommand*\pandocbounded[1]{% scales image to fit in text height/width
  \sbox\pandoc@box{#1}%
  \Gscale@div\@tempa{\textheight}{\dimexpr\ht\pandoc@box+\dp\pandoc@box\relax}%
  \Gscale@div\@tempb{\linewidth}{\wd\pandoc@box}%
  \ifdim\@tempb\p@<\@tempa\p@\let\@tempa\@tempb\fi% select the smaller of both
  \ifdim\@tempa\p@<\p@\scalebox{\@tempa}{\usebox\pandoc@box}%
  \else\usebox{\pandoc@box}%
  \fi%
}
% Set default figure placement to htbp
\def\fps@figure{htbp}
\makeatother
\setlength{\emergencystretch}{3em} % prevent overfull lines
\providecommand{\tightlist}{%
  \setlength{\itemsep}{0pt}\setlength{\parskip}{0pt}}
\usepackage{bookmark}
\IfFileExists{xurl.sty}{\usepackage{xurl}}{} % add URL line breaks if available
\urlstyle{same}
% fallback for those not using the hyperref driver hyperxmp:
\makeatletter
\@ifundefined{xmpquote}{\newcommand{\xmpquote}[1]{#1}}{}
\makeatother
\hypersetup{
  hidelinks,
  pdfcreator={LaTeX via pandoc}}

\author{}
\date{}

\begin{document}

\section{Ekspansi Industri
Ekstraktif}\label{ekspansi-industri-ekstraktif}

\emph{Analisis spasiotemporal pertumbuhan industri ekstraktif dan
pengolahan nikel serta dampaknya terhadap daya dukung dan daya tampung
lingkungan di Pulau Sulawesi.}

\subsection{Ekspansi Industri Ekstraktif: 778 Unit Smelter dan
Ketergantungan Energi Fosil Off-Grid di
Sulawesi}\label{ekspansi-industri-ekstraktif-778-unit-smelter-dan-ketergantungan-energi-fosil-off-grid-di-sulawesi}

Dinamika pembangunan di Pulau Sulawesi periode 2014--2024 ditandai oleh
akselerasi industri berbasis komoditas alam. Kebijakan hilirisasi nikel
mendorong penerbitan \textbf{574 Izin Usaha Pertambangan (IUP) baru}
dengan total luas konsesi mencapai \textbf{819,452 Hektar}.
Pengoperasian \textbf{778 unit fasilitas pemurnian (smelter)} didukung
oleh kapasitas \textbf{9,825 MW PLTU Captive} (pembangkit listrik batu
bara \emph{off-grid}), yang meningkatkan intensitas emisi karbon pada
zona-zona industri pesisir.

Secara bersamaan, kucuran realisasi Penanaman Modal Dalam Negeri (PMDN)
yang mencapai \textbf{218 Triliun Rupiah} berbanding lurus dengan
akumulasi konversi tutupan hutan sebesar \textbf{2,107,041 Hektar} untuk
aktivitas pertambangan dan perkebunan. Data ini mengindikasikan bahwa
pertumbuhan indikator makroekonomi berjalan seiring dengan peningkatan
beban terhadap daya dukung dan daya tampung lingkungan hidup.

\subsubsection{Metrik Ekstraktif}\label{metrik-ekstraktif}

{\def\LTcaptype{none} % do not increment counter
\begin{longtable}[]{@{}
  >{\raggedright\arraybackslash}p{(\linewidth - 4\tabcolsep) * \real{0.3333}}
  >{\raggedright\arraybackslash}p{(\linewidth - 4\tabcolsep) * \real{0.3333}}
  >{\raggedright\arraybackslash}p{(\linewidth - 4\tabcolsep) * \real{0.3333}}@{}}
\toprule\noalign{}
\begin{minipage}[b]{\linewidth}\raggedright
Indikator
\end{minipage} & \begin{minipage}[b]{\linewidth}\raggedright
Nilai
\end{minipage} & \begin{minipage}[b]{\linewidth}\raggedright
Deskripsi
\end{minipage} \\
\midrule\noalign{}
\endhead
\bottomrule\noalign{}
\endlastfoot
\textbf{Total Izin Baru (2014-2024)} & \textbf{574 IUP} & Penambahan
jumlah IUP di Pulau Sulawesi dalam 1 dekade terakhir. \\
\textbf{Total Luas Konsesi Baru} & \textbf{819,452 Ha} & Akumulasi luas
daratan dan pesisir yang diserahkan sejak 2014. \\
\textbf{Kapasitas PLTU Captive Aktif} & \textbf{9,825 MW} & Beban energi
kotor off-grid untuk menyokong pabrik peleburan. \\
\textbf{Jumlah Fasilitas Smelter} & \textbf{778 Unit} & Total fasilitas
pengolahan nikel yang memonopoli zona industri pesisir. \\
\textbf{Luas Deforestasi Komoditas} & \textbf{2,107,041 Ha} & Area hutan
alam yang musnah akibat tambang dan perkebunan. \\
\textbf{Investasi PMDN (2016-2024)} & \textbf{218 Triliun Rp} & Aliran
modal domestik yang dikucurkan. \\
\end{longtable}
}

\begin{center}\rule{0.5\linewidth}{0.5pt}\end{center}

\subsection{1.1 Konteks Makro: Breakdown PDRB per
Komoditas}\label{konteks-makro-breakdown-pdrb-per-komoditas}

Grafik di bawah ini menyederhanakan 17 sektor PDRB menjadi \textbf{3
klasifikasi makro advokatif} berdasarkan \emph{Legal Supply-Chain
Approach} (Metodologi CELIOS/ECC).

\subsubsection{1.1.1 Dominasi Ekstraktif vs Ekonomi Akar Rumput
(2016-2024)}\label{dominasi-ekstraktif-vs-ekonomi-akar-rumput-2016-2024}

Grafik di bawah ini menyederhanakan 17 sektor PDRB menjadi \textbf{3
klasifikasi makro advokatif} berdasarkan \emph{Legal Supply-Chain
Approach} (Metodologi CELIOS/ECC).

\begin{itemize}
\tightlist
\item
  \textbf{Ekstraktif} = Kat. B (Pertambangan) + Kat. C (Industri
  Pengolahan/Smelter) + Kat. D (Listrik/PLTU Captive) --- digabung
  berdasarkan mandat wajib UU Minerba Ps. 102-103 \& Perpres 112/2022.
\item
  \textbf{Ekonomi Akar Rumput} = Kat. A (Pertanian, Kehutanan \&
  Perikanan) --- sektor terbarukan penyerap tenaga kerja lokal terbesar.
\item
  \textbf{Sektor Jasa \& Lainnya} = 13 sektor E-U sisanya.
\end{itemize}

\begin{figure}
\centering
\pandocbounded{\includegraphics[keepaspectratio,alt={1.1.1}]{visuals/chart_1_1_1.png}}
\caption{1.1.1}
\end{figure}

\emph{Metodologi: Legal Supply-Chain Approach --- Kat B+C+D = Ekstraktif
(UU Minerba Ps.102-103; Perpres 112/2022 Ps.3 Ay.4)}

\subsubsection{1.1.2 Pemusatan Sektor Ekstraktif di Kabupaten
se-Sulawesi
Tengah}\label{pemusatan-sektor-ekstraktif-di-kabupaten-se-sulawesi-tengah}

Visualisasi di bawah membandingkan komposisi ketiga sektor advokatif di
seluruh 13 kabupaten/kota se-Sulawesi Tengah pada tahun terbaru.

\begin{figure}
\centering
\pandocbounded{\includegraphics[keepaspectratio,alt={1.1.2}]{visuals/chart_1_1_2.png}}
\caption{1.1.2}
\end{figure}

\subsubsection{1.1.3 Perbandingan Distribusi 17 Sektor Komoditas per
Provinsi (Small Multiples, Tahun
Terbaru)}\label{perbandingan-distribusi-17-sektor-komoditas-per-provinsi-small-multiples-tahun-terbaru}

Visualisasi Small Multiples ini membandingkan komposisi 17 sektor
komoditas secara terpisah di tiap provinsi. Sektor diurutkan dari
penyumbang terbesar (atas) hingga terkecil (bawah). Skala sumbu X
konsisten untuk memvalidasi perbandingan lintas provinsi.

\begin{figure}
\centering
\pandocbounded{\includegraphics[keepaspectratio,alt={1.1.3}]{visuals/chart_1_1_3.png}}
\caption{1.1.3}
\end{figure}

\begin{center}\rule{0.5\linewidth}{0.5pt}\end{center}

\subsection{1.2 Konsentrasi Kawasan Industri \& PLTU
Captive}\label{konsentrasi-kawasan-industri-pltu-captive}

Intensifikasi industri pengolahan nikel di Sulawesi berpusat pada
fasilitas mega-smelter. Pengoperasian \textbf{778 fasilitas smelter}
didukung oleh kapasitas energi batu bara \textbf{9,825 MW dari PLTU
Captive}. Berbeda dengan sistem kelistrikan umum PLN, pembangkit ini
dikembangkan secara internal untuk menyokong operasi kawasan industri.

Berikut adalah \textbf{temuan konsentrasi spasial} berdasarkan data
agregat:

\textbf{Pemusatan Spasial Fasilitas Smelter (Bar Chart):} Data
menunjukkan bahwa \textbf{77\% dari total fasilitas (344 unit di
Sulawesi Tengah dan 262 unit di Sulawesi Tenggara)} terkonsentrasi di
dua provinsi tersebut. Pola ini mengonfirmasi adanya pemusatan beban
ekologis dan emisi pada zona sentra pemurnian nikel.

Korelasi antara pembangunan kawasan industri dan perubahan tutupan lahan
diuji menggunakan \textbf{Crosstabulation (Tabulasi Silang)} pada bagian
bawah sub-bab ini.

\begin{figure}
\centering
\pandocbounded{\includegraphics[keepaspectratio,alt={1.2}]{visuals/chart_1_2.png}}
\caption{1.2}
\end{figure}

\textbf{Fakta Data:} Sebesar 78\% dari total 778 fasilitas smelter
terkonsentrasi di Sulawesi Tengah \& Sulawesi Tenggara, menunjukkan
adanya pemusatan beban lingkungan di wilayah sentra tersebut.

\subsubsection{Pembuktian Statistik: Ekspansi PLTU Captive vs
Deforestasi}\label{pembuktian-statistik-ekspansi-pltu-captive-vs-deforestasi}

Untuk menguji apakah keberadaan PLTU \emph{Captive} berkorelasi secara
spasial dan temporal dengan laju deforestasi, kita menggunakan tabel
crosstab pada level observasi \textbf{Provinsi-Tahun}. Mengingat
ekspansi PLTU sangat terpusat pada tahun dan provinsi tertentu
(menghasilkan banyak nilai nol pada panel), klasifikasi ``Tinggi''
diartikan sebagai \emph{ada penambahan kapasitas (\textgreater0)}, dan
``Rendah'' sebagai \emph{tidak ada penambahan (=0)}.

\textbf{Case Processing Summary}

{\def\LTcaptype{none} % do not increment counter
\begin{longtable}[]{@{}lll@{}}
\toprule\noalign{}
Keterangan & N & Persen \\
\midrule\noalign{}
\endhead
\bottomrule\noalign{}
\endlastfoot
Valid & 60 & 100.0\% \\
Total & 60 & 100.0\% \\
\end{longtable}
}

\textbf{Chi-Square Tests}

{\def\LTcaptype{none} % do not increment counter
\begin{longtable}[]{@{}llll@{}}
\toprule\noalign{}
Uji & Value & df & Asymp. Sig. (2-sided) \\
\midrule\noalign{}
\endhead
\bottomrule\noalign{}
\endlastfoot
Pearson Chi-Square & 0 & 0 & 1 \\
\end{longtable}
}

\textbf{Result: TIDAK SIGNIFIKAN} (P-Value = 1, Odds Ratio = 0)

\emph{Meski data tahunan agregat menunjukkan tidak signifikan
(kemungkinan karena konsentrasi PLTU hanya terjadi di segelintir tahun
dan lokasi seperti Morowali), hal ini bukan berarti PLTU ramah
lingkungan. Sebaliknya, efek rusak dari sebuah PLTU bersifat permanen
dan lintas-batas (spillover) yang mencemari wilayah di luar lokasi
spesifik pendiriannya.}

\textbf{Interpretasi Spasial Industri:} Kawasan industri pengolahan
terkonsentrasi di area pesisir secara signifikan. Pertumbuhan PLTU
Captive mengindikasikan tingginya ketergantungan pada energi berbasis
batu bara untuk mendukung kebutuhan energi fasilitas pemurnian di
Sulawesi Tengah dan Sulawesi Tenggara.

\subsubsection{Ringkasan Eksekutif Seluruh Skenario
Crosstab}\label{ringkasan-eksekutif-seluruh-skenario-crosstab}

Tabel di bawah ini merangkum hasil pengujian statistik (Chi-Square)
untuk semua kemungkinan kombinasi indikator antara penambahan PLTU
Captive dan Dampak Ekologis pada panel data 1 dekade.

{\def\LTcaptype{none} % do not increment counter
\begin{longtable}[]{@{}
  >{\raggedright\arraybackslash}p{(\linewidth - 10\tabcolsep) * \real{0.1667}}
  >{\raggedright\arraybackslash}p{(\linewidth - 10\tabcolsep) * \real{0.1667}}
  >{\raggedright\arraybackslash}p{(\linewidth - 10\tabcolsep) * \real{0.1667}}
  >{\raggedright\arraybackslash}p{(\linewidth - 10\tabcolsep) * \real{0.1667}}
  >{\raggedright\arraybackslash}p{(\linewidth - 10\tabcolsep) * \real{0.1667}}
  >{\raggedright\arraybackslash}p{(\linewidth - 10\tabcolsep) * \real{0.1667}}@{}}
\toprule\noalign{}
\begin{minipage}[b]{\linewidth}\raggedright
Variabel Independen (X)
\end{minipage} & \begin{minipage}[b]{\linewidth}\raggedright
Variabel Dependen (Y)
\end{minipage} & \begin{minipage}[b]{\linewidth}\raggedright
Chi-Square
\end{minipage} & \begin{minipage}[b]{\linewidth}\raggedright
P-Value
\end{minipage} & \begin{minipage}[b]{\linewidth}\raggedright
Odds Ratio
\end{minipage} & \begin{minipage}[b]{\linewidth}\raggedright
Kesimpulan
\end{minipage} \\
\midrule\noalign{}
\endhead
\bottomrule\noalign{}
\endlastfoot
Kapasitas Aktif PLTU Kumulatif (MW) & Total Deforestasi Alam (Hektar) &
0.000 & 1.000 & 0.00 & ❌ TIDAK SIGNIFIKAN \\
Kapasitas Aktif PLTU Kumulatif (MW) & Deforestasi Komoditas
Tambang/Sawit (Hektar) & 0.000 & 1.000 & 0.00 & ❌ TIDAK SIGNIFIKAN \\
\end{longtable}
}

Dari \textbf{2 skenario pengujian}, seluruhnya menunjukkan status
\textbf{TIDAK SIGNIFIKAN}. Meski demikian, ini bukan berarti PLTU ramah
lingkungan---efek destruktifnya bersifat \emph{spillover} yang mencemari
wilayah bahkan di luar lokasi spesifik pendiriannya.

\begin{center}\rule{0.5\linewidth}{0.5pt}\end{center}

\subsection{1.3 Tren Pertumbuhan Izin Tambang Baru \& Uji
Signifikansi}\label{tren-pertumbuhan-izin-tambang-baru-uji-signifikansi}

Pola perizinan pertambangan di Pulau Sulawesi selama satu dekade
terakhir menunjukkan peningkatan alokasi ruang yang signifikan.
Berdasarkan data agregat \emph{Minerbaone}, tercatat \textbf{574 Izin
Usaha Pertambangan (IUP) baru} sepanjang 2014-2024, dengan total luas
konsesi mencapai \textbf{819,452 Hektar}.

Berdasarkan analisis tren time-series pada grafik \textbf{``Penerbitan
Izin Tambang''} di bawah, penerbitan izin pada periode awal (2014)
tercatat sebanyak \textbf{26 IUP}. Peningkatan signifikan terjadi pada
periode 2022--2024, di mana penerbitan meningkat dari \textbf{56 IUP di
tahun 2022} menjadi \textbf{149 IUP pada 2023}, dan mencapai \textbf{194
IUP baru pada 2024}.

Anotasi pada grafik mencatat kenaikan sebesar \textbf{246\% pada periode
2022--2024}. Data ini mengindikasikan perlunya evaluasi terhadap
instrumen pengendalian perizinan dan tata ruang. Distribusi perizinan
tertinggi berada di Sulawesi Tengah dan Sulawesi Tenggara, yang selaras
dengan kawasan pengembangan industri pemurnian nikel.

Uji \textbf{Crosstabulation} pada bagian bawah mengukur hubungan antara
laju penerbitan perizinan dan indikator deforestasi di wilayah tersebut.

\begin{figure}
\centering
\pandocbounded{\includegraphics[keepaspectratio,alt={1.3}]{visuals/chart_1_3.png}}
\caption{1.3}
\end{figure}

\textbf{Interpretasi Sektoral:} Peningkatan penerbitan IUP di kawasan
timur Sulawesi berbanding lurus dengan perluasan area konversi hutan.
Pola perizinan ini menunjukkan pentingnya penerapan instrumen tata ruang
dan evaluasi lingkungan secara ketat.

\subsubsection{Pembuktian Statistik: Intensitas Ekspansi vs
Deforestasi}\label{pembuktian-statistik-intensitas-ekspansi-vs-deforestasi}

Hipotesis utama narasi ini adalah bahwa \textbf{lonjakan ekspansi
ekstraktif} berbanding lurus dengan \textbf{kebangkrutan ekologis}
(deforestasi). Untuk mengujinya secara statistik di tengah keterbatasan
jumlah provinsi di Sulawesi (N=6), tabel crosstab dan uji Chi-Square di
bawah menggunakan unit observasi \textbf{Provinsi-Tahun} (6 provinsi x
10 tahun = 60 sampel panel). Setiap observasi diklasifikasikan menjadi
``Tinggi'' atau ``Rendah'' berdasarkan nilai \textbf{Median panel} dari
indikator yang dipilih.

\textbf{Chi-Square Tests (Jumlah Izin Baru IUP * Deforestasi Komoditas)}

{\def\LTcaptype{none} % do not increment counter
\begin{longtable}[]{@{}llll@{}}
\toprule\noalign{}
Uji & Value & df & Asymp. Sig. (2-sided) \\
\midrule\noalign{}
\endhead
\bottomrule\noalign{}
\endlastfoot
Pearson Chi-Square & 1.068 & 1 & 0.301 \\
N of Valid Cases & 60 & & \\
\end{longtable}
}

\textbf{Result: TIDAK SIGNIFIKAN} (P-Value = 0.3014, Odds Ratio = 1.962)

\emph{Secara agregat, hubungan antara Jumlah Izin Baru dan Deforestasi
Komoditas tidak signifikan secara statistik (P ≥ 0.05). Ini
mengindikasikan bahwa deforestasi terjadi sangat masif di seluruh panel
waktu dan ruang secara merata. Krisis tata kelola dan deforestasi telah
menyebar ke seluruh wilayah, sehingga lonjakan izin di tahun tertentu
tidak lagi menjadi prediktor tunggal atas kebangkrutan ekologis yang
sudah sistemik.}

\subsubsection{Ringkasan Eksekutif Seluruh Skenario
Crosstab}\label{ringkasan-eksekutif-seluruh-skenario-crosstab-1}

Tabel di bawah ini merangkum hasil pengujian statistik (Chi-Square)
untuk semua kemungkinan kombinasi antara indikator Ekspansi (X) dan
Dampak Ekologis (Y) pada panel data yang sama.

{\def\LTcaptype{none} % do not increment counter
\begin{longtable}[]{@{}
  >{\raggedright\arraybackslash}p{(\linewidth - 10\tabcolsep) * \real{0.1667}}
  >{\raggedright\arraybackslash}p{(\linewidth - 10\tabcolsep) * \real{0.1667}}
  >{\raggedright\arraybackslash}p{(\linewidth - 10\tabcolsep) * \real{0.1667}}
  >{\raggedright\arraybackslash}p{(\linewidth - 10\tabcolsep) * \real{0.1667}}
  >{\raggedright\arraybackslash}p{(\linewidth - 10\tabcolsep) * \real{0.1667}}
  >{\raggedright\arraybackslash}p{(\linewidth - 10\tabcolsep) * \real{0.1667}}@{}}
\toprule\noalign{}
\begin{minipage}[b]{\linewidth}\raggedright
Variabel Independen (X)
\end{minipage} & \begin{minipage}[b]{\linewidth}\raggedright
Variabel Dependen (Y)
\end{minipage} & \begin{minipage}[b]{\linewidth}\raggedright
Chi-Square
\end{minipage} & \begin{minipage}[b]{\linewidth}\raggedright
P-Value
\end{minipage} & \begin{minipage}[b]{\linewidth}\raggedright
Odds Ratio
\end{minipage} & \begin{minipage}[b]{\linewidth}\raggedright
Kesimpulan
\end{minipage} \\
\midrule\noalign{}
\endhead
\bottomrule\noalign{}
\endlastfoot
Jumlah Izin Baru (IUP) & Total Deforestasi Alam (Hektar) & 17.239 &
0.000 & 13.75 & ✅ SIGNIFIKAN \\
Jumlah Izin Baru (IUP) & Deforestasi Komoditas Tambang/Sawit (Hektar) &
1.077 & 0.299 & 1.97 & ❌ TIDAK SIGNIFIKAN \\
Luas Konsesi Baru (Hektar) & Total Deforestasi Alam (Hektar) & 11.267 &
0.001 & 7.56 & ✅ SIGNIFIKAN \\
Luas Konsesi Baru (Hektar) & Deforestasi Komoditas Tambang/Sawit
(Hektar) & 3.267 & 0.071 & 2.98 & ❌ TIDAK SIGNIFIKAN \\
\end{longtable}
}

Dari \textbf{4 skenario pengujian}, terdapat \textbf{2 skenario yang
terbukti SIGNIFIKAN}. Tingginya \emph{Odds Ratio} pada skenario
signifikan menegaskan bahwa setiap kali kran perizinan atau luas konsesi
diperlebar, risiko terjadinya deforestasi melonjak berkali-kali lipat.

\begin{center}\rule{0.5\linewidth}{0.5pt}\end{center}

\subsection{1.4 Analisis Realisasi Investasi PMDN dan Dampak Terhadap
Tutupan
Hutan}\label{analisis-realisasi-investasi-pmdn-dan-dampak-terhadap-tutupan-hutan}

Grafik di bawah ini memetakan dinamika penerbitan konsesi tambang baru
dan dampaknya terhadap tutupan hutan. Pada tahun 2016, luas konsesi
tambang baru yang diterbitkan di Sulawesi mencakup \textbf{12,515
Hektar}, dan meningkat signifikan hingga mencapai \textbf{68,556 Hektar}
pada tahun 2023. Pada periode yang sama, angka deforestasi komoditas
mencatatkan luasan sebesar \textbf{216,887 Hektar}.

Data ini mengindikasikan bahwa akselerasi penerbitan konsesi berbanding
lurus dengan laju konversi hutan alam (akumulasi deforestasi sebesar
\textbf{2,107,041 Hektar}). Hal ini menegaskan pentingnya pertimbangan
daya dukung ekologis dalam setiap kebijakan alokasi konsesi
pertambangan.

\begin{figure}
\centering
\pandocbounded{\includegraphics[keepaspectratio,alt={1.4}]{visuals/chart_1_4.png}}
\caption{1.4}
\end{figure}

\textbf{Interpretasi Spasial:} Perbandingan grafik batang di atas
menunjukkan bahwa tingkat alokasi konsesi di Daerah Sentra Tambang
(Morowali \& Konawe) jauh lebih tinggi dibanding wilayah non-sentra,
yang berdampak langsung pada konsentrasi perubahan tutupan hutan.

\subsubsection{Pembuktian Statistik: Arus Investasi PMDN vs
Deforestasi}\label{pembuktian-statistik-arus-investasi-pmdn-vs-deforestasi}

\textbf{Chi-Square Tests (Realisasi Investasi PMDN * Deforestasi
Komoditas)}

{\def\LTcaptype{none} % do not increment counter
\begin{longtable}[]{@{}llll@{}}
\toprule\noalign{}
Uji & Value & df & Asymp. Sig. (2-sided) \\
\midrule\noalign{}
\endhead
\bottomrule\noalign{}
\endlastfoot
Pearson Chi-Square & 7.042 & 1 & 0.008 \\
N of Valid Cases & 96 & & \\
\end{longtable}
}

\textbf{Result: SIGNIFIKAN} (P-Value = 0.008, Odds Ratio = 3.325)

\emph{Terdapat bukti statistik yang sah bahwa arus masuk modal
(Investasi PMDN) secara langsung dan sistematis mendorong ekskalasi
deforestasi di wilayah Sulawesi (OR: 3.325). Investasi ini bukanlah
katalisator ekonomi hijau, melainkan injeksi modal untuk ekstraksi
lahan.}

\subsubsection{Ringkasan Eksekutif Seluruh Skenario Crosstab (Investasi
vs
Deforestasi)}\label{ringkasan-eksekutif-seluruh-skenario-crosstab-investasi-vs-deforestasi}

Tabel di bawah ini merangkum hasil pengujian statistik (Chi-Square)
untuk semua kemungkinan kombinasi antara Realisasi Investasi PMDN dan
Dampak Ekologis pada panel data.

{\def\LTcaptype{none} % do not increment counter
\begin{longtable}[]{@{}
  >{\raggedright\arraybackslash}p{(\linewidth - 10\tabcolsep) * \real{0.1667}}
  >{\raggedright\arraybackslash}p{(\linewidth - 10\tabcolsep) * \real{0.1667}}
  >{\raggedright\arraybackslash}p{(\linewidth - 10\tabcolsep) * \real{0.1667}}
  >{\raggedright\arraybackslash}p{(\linewidth - 10\tabcolsep) * \real{0.1667}}
  >{\raggedright\arraybackslash}p{(\linewidth - 10\tabcolsep) * \real{0.1667}}
  >{\raggedright\arraybackslash}p{(\linewidth - 10\tabcolsep) * \real{0.1667}}@{}}
\toprule\noalign{}
\begin{minipage}[b]{\linewidth}\raggedright
Variabel Independen (X)
\end{minipage} & \begin{minipage}[b]{\linewidth}\raggedright
Variabel Dependen (Y)
\end{minipage} & \begin{minipage}[b]{\linewidth}\raggedright
Chi-Square
\end{minipage} & \begin{minipage}[b]{\linewidth}\raggedright
P-Value
\end{minipage} & \begin{minipage}[b]{\linewidth}\raggedright
Odds Ratio
\end{minipage} & \begin{minipage}[b]{\linewidth}\raggedright
Kesimpulan
\end{minipage} \\
\midrule\noalign{}
\endhead
\bottomrule\noalign{}
\endlastfoot
Realisasi Investasi PMDN (Juta Rp) & Total Deforestasi Alam (Hektar) &
3.375 & 0.066 & 2.33 & ❌ TIDAK SIGNIFIKAN \\
Realisasi Investasi PMDN (Juta Rp) & Deforestasi Komoditas Tambang/Sawit
(Hektar) & 7.042 & 0.008 & 3.33 & ✅ SIGNIFIKAN \\
\end{longtable}
}

Dari \textbf{2 skenario pengujian}, terdapat \textbf{1 skenario yang
terbukti SIGNIFIKAN}. Derasnya arus modal (PMDN) bukan indikator
keberhasilan ekonomi yang inklusif, melainkan sekadar dana segar untuk
membiayai penghancuran hutan skala raksasa.

\begin{center}\rule{0.5\linewidth}{0.5pt}\end{center}

\subsection{1.5 Pelabuhan Ekspor: Ke Mana Nikel Sulawesi
Dikirim?}\label{pelabuhan-ekspor-ke-mana-nikel-sulawesi-dikirim}

Ekspansi nikel di Sulawesi tidak berhenti pada izin dan pabrik smelter.
Di setiap lokasi industri nikel besar, berdiri \textbf{pelabuhan atau
dermaga} yang menghubungkan pabrik langsung ke kapal-kapal pengangkut
menuju China dan pasar global. Dari 6 lokasi utama yang ditelusuri,
\textbf{seluruhnya terbukti memiliki} pelabuhan atau dermaga ekspor, dan
\textbf{4 dari 6} mendapat label Proyek Strategis Nasional (PSN) dari
pemerintah.

{\def\LTcaptype{none} % do not increment counter
\begin{longtable}[]{@{}ll@{}}
\toprule\noalign{}
Indikator & Nilai \\
\midrule\noalign{}
\endhead
\bottomrule\noalign{}
\endlastfoot
\textbf{Pelabuhan Nikel Terkonfirmasi} & \textbf{6} Lokasi \\
\textbf{Berlabel Proyek Strategis Nasional} & \textbf{4 / 6} \\
\textbf{Kapasitas Pelabuhan Terbesar} & \textbf{50.000 ton} (GNI
Petasia) \\
\end{longtable}
}

\emph{Sumber: Situs perusahaan, dokumen pemerintah, media (25 sumber
OSINT). File: sulawesi\_logistik\_simpul\_nikel.csv}

\begin{center}\rule{0.5\linewidth}{0.5pt}\end{center}

\subsection{1.6 Peta Jalur Distribusi Logistik Nikel
Sulawesi}\label{peta-jalur-distribusi-logistik-nikel-sulawesi}

Peta rute logistik maritim ini mengilustrasikan realitas geopolitik dari
ambisi hilirisasi nikel di Sulawesi. Alih-alih membangun kemandirian
industri manufaktur nasional, data pergerakan kapal dan desain pelabuhan
menunjukkan \textbf{ketergantungan absolut pada rantai pasok asing}.

{\def\LTcaptype{none} % do not increment counter
\begin{longtable}[]{@{}
  >{\raggedright\arraybackslash}p{(\linewidth - 6\tabcolsep) * \real{0.2500}}
  >{\raggedright\arraybackslash}p{(\linewidth - 6\tabcolsep) * \real{0.2500}}
  >{\raggedright\arraybackslash}p{(\linewidth - 6\tabcolsep) * \real{0.2500}}
  >{\raggedright\arraybackslash}p{(\linewidth - 6\tabcolsep) * \real{0.2500}}@{}}
\toprule\noalign{}
\begin{minipage}[b]{\linewidth}\raggedright
Nama Smelter
\end{minipage} & \begin{minipage}[b]{\linewidth}\raggedright
Asal (Lon, Lat)
\end{minipage} & \begin{minipage}[b]{\linewidth}\raggedright
Tujuan
\end{minipage} & \begin{minipage}[b]{\linewidth}\raggedright
Komoditas
\end{minipage} \\
\midrule\noalign{}
\endhead
\bottomrule\noalign{}
\endlastfoot
\textbf{IMIP} & Morowali (122.15, -2.82) & China & NPI/Feronikel \\
\textbf{GNI} & Morowali Utara (121.32, -1.91) & China & NPI \\
\textbf{VDNI} & Konawe (122.42, -3.83) & China & Feronikel \& Stainless
Steel \\
\textbf{OSS} & Konawe (122.48, -3.80) & China & Stainless Steel \\
\textbf{ANTAM} & Kolaka (121.60, -4.18) & Jepang/Korea & Feronikel \\
\textbf{PT Vale} & Luwu Timur (121.34, -2.56) & Jepang/Korea & Nickel in
Matte \\
\end{longtable}
}

\begin{figure}
\centering
\pandocbounded{\includegraphics[keepaspectratio,alt={Peta Jalur Distribusi Logistik Nikel Sulawesi}]{visuals/chart_1_6_peta.png}}
\caption{Peta Jalur Distribusi Logistik Nikel Sulawesi}
\end{figure}

\textbf{Ketergantungan Struktural Rantai Pasok:} - \textbf{Dominasi
Ekspor ke China:} Tiga raksasa kawasan industri baru (IMIP, GNI,
VDNI/OSS) yang menikmati fasilitas kemudahan Proyek Strategis Nasional
(PSN) mengirimkan hampir seluruh \emph{output} barang setengah jadi
(NPI, Feronikel, Matte) langsung ke sentra industri di China Timur dan
Selatan. - \textbf{Absennya Interkoneksi Domestik:} Sangat minim jalur
distribusi logistik yang menghubungkan kawasan smelter raksasa ini
dengan pusat industri manufaktur di dalam negeri (seperti di Pulau
Jawa). Hal ini mengonfirmasi temuan bahwa Sulawesi saat ini lebih
difungsikan murni sebagai \emph{extractive feeder} (daerah penyuplai
ekstraktif) bagi mesin industrialisasi negara lain, bukan sebagai
fondasi terintegrasi untuk ekosistem mobil listrik domestik. -
\textbf{Pergeseran Geopolitik:} Sementara pemain lama seperti PT Vale
dan ANTAM memiliki rute pasokan yang mapan ke pasar otomotif tradisional
di Jepang dan Korea Selatan, dominasi logistik dan tonase kini telah
bergeser drastis seiring dengan peningkatan signifikan pembangunan
smelter baru yang terintegrasi langsung dengan pasar China.

\end{document}
